% GENERATED FILE. Edit canonical lessons, then run npm run generate:books.
% canonical-source-hash: fnv1a64:fa5a58557d55eb7a
% canonical-lessons: RU-C527-serial, RU-C527-epizod, RU-C527-subtitry, RU-C527-pereklyuchat, RU-C527-ekranizatsiya

\chapter{Series and Subtitles}
\label{ch:ru-series-and-subtitles}
\hlchaptermodality{527}

\begin{chapteropening}
\textbf{By the end of this chapter:} \emph{I can talk about watching series and adaptations.}

Complete the last lesson of chapter 527: I can talk about watching series and adaptations.
\end{chapteropening}

\section[\texorpdfstring{\ru{сериал} (seriál)}{сериал (seriál)}]{\ru{сериал} (seriál) — a TV series}

\label{lesson:RU-C527-serial}



\begin{quote}
Before the new one: say the Russian for a publishing house, then the Russian for to reread.
\end{quote}

\subsection*{You'll want to know: \ru{сериал}}
\textbf{\ru{сериал}} (\emph{seriál}) — \textquotedblleft{}a TV series\textquotedblright{}.

\begin{quote}
\ru{Мы} \ru{смотрим} \ru{этот} \ru{сериал} \ru{уже} \ru{год}.
\end{quote}

\begin{quote}
\emph{My smótrim étot seriál uzhé god.}
\end{quote}

\begin{quote}
We've been watching this series for a year.
\end{quote}

Masculine.

\begin{cousinweb}
\emph{\ru{сериал}} is usually traced to English \emph{serial}; EtymDB has no entry.
\end{cousinweb}

\begin{grammarlens}[title={What you've built}]
One more word for telling, explaining and describing at B1.
\end{grammarlens}

\subsection*{Your turn}
\begin{itemize}
\raggedright
  \item \emph{Say it:} \emph{seriál}
  \item \emph{Say it:} \emph{seriál}, once more
  \item \emph{Say it:} a sentence of your own with \emph{seriál}
  \item \emph{Recall:} say the Russian for a publishing house, then the Russian for to reread, then say \emph{seriál} again
\end{itemize}

\begin{tcolorbox}[breakable,colback=teal!4,colframe=teal!35!black,title={Before you move on}]
What does \ru{сериал} mean? (\textquotedblleft{}A TV series\textquotedblright{}.) Say it once more.
\end{tcolorbox}
\section[\texorpdfstring{\ru{эпизод} (epizód)}{эпизод (epizód)}]{\ru{эпизод} (epizód) — an episode}

\label{lesson:RU-C527-epizod}



\begin{quote}
Before the new one: say the Russian for to reread, then the Russian for a TV series.
\end{quote}

\subsection*{You'll want to know: \ru{эпизод}}
\textbf{\ru{эпизод}} (\emph{epizód}) — \textquotedblleft{}an episode\textquotedblright{}.

\begin{quote}
\ru{Я} \ru{посмотрел} \ru{три} \ru{эпизода} \ru{подряд}.
\end{quote}

\begin{quote}
\emph{Ya posmotrél tri epizóda podryád.}
\end{quote}

\begin{quote}
I watched three episodes in a row.
\end{quote}

Masculine. Also a short scene or incident.

\begin{cousinweb}
EtymDB lists \emph{\ru{эпизод}} without a source; it is usually traced through French \emph{épisode} to Greek \emph{epeisódion} 'an added scene'.
\end{cousinweb}

\begin{grammarlens}[title={What you've built}]
One more word for telling, explaining and describing at B1.
\end{grammarlens}

\subsection*{Your turn}
\begin{itemize}
\raggedright
  \item \emph{Say it:} \emph{epizód}
  \item \emph{Say it:} \emph{epizód}, once more
  \item \emph{Say it:} a sentence of your own with \emph{epizód}
  \item \emph{Recall:} say the Russian for to reread, then the Russian for a TV series, then say \emph{epizód} again
\end{itemize}

\begin{tcolorbox}[breakable,colback=teal!4,colframe=teal!35!black,title={Before you move on}]
What does \ru{эпизод} mean? (\textquotedblleft{}An episode\textquotedblright{}.) Say it once more.
\end{tcolorbox}
\section[\texorpdfstring{\ru{субтитры} (subtítry)}{субтитры (subtítry)}]{\ru{субтитры} (subtítry) — subtitles}

\label{lesson:RU-C527-subtitry}



\begin{quote}
Before the new one: say the Russian for a TV series, then the Russian for an episode.
\end{quote}

\subsection*{You'll want to know: \ru{субтитры}}
\textbf{\ru{субтитры}} (\emph{subtítry}) — \textquotedblleft{}subtitles\textquotedblright{}.

\begin{quote}
\ru{Я} \ru{смотрю} \ru{фильмы} \ru{с} \ru{субтитрами}.
\end{quote}

\begin{quote}
\emph{Ya smotryú fílmy s subtítrami.}
\end{quote}

\begin{quote}
I watch films with subtitles.
\end{quote}

Plural.

\begin{cousinweb}
\emph{\ru{субтитры}} is built on \emph{\ru{суб}-} 'under' and \emph{\ru{титр}} 'a caption', usually traced through French \emph{titre} to Latin \emph{titulus} 'a label'; EtymDB has no entry.
\end{cousinweb}

\begin{grammarlens}[title={What you've built}]
One more word for telling, explaining and describing at B1.
\end{grammarlens}

\subsection*{Your turn}
\begin{itemize}
\raggedright
  \item \emph{Say it:} \emph{subtítry}
  \item \emph{Say it:} \emph{subtítry}, once more
  \item \emph{Say it:} a sentence of your own with \emph{subtítry}
  \item \emph{Recall:} say the Russian for a TV series, then the Russian for an episode, then say \emph{subtítry} again
\end{itemize}

\begin{tcolorbox}[breakable,colback=teal!4,colframe=teal!35!black,title={Before you move on}]
What does \ru{субтитры} mean? (\textquotedblleft{}Subtitles\textquotedblright{}.) Say it once more.
\end{tcolorbox}
\section[\texorpdfstring{\ru{переключать} (pereklyuchát)}{переключать (pereklyuchát)}]{\ru{переключать} (pereklyuchát) — to switch (channels)}

\label{lesson:RU-C527-pereklyuchat}



\begin{quote}
Before the new one: say the Russian for an episode, then the Russian for subtitles.
\end{quote}

\subsection*{You'll want to know: \ru{переключать}}
\textbf{\ru{переключать}} (\emph{pereklyuchát}) — \textquotedblleft{}to switch (channels)\textquotedblright{}.

\begin{quote}
\ru{Папа} \ru{всё} \ru{время} \ru{переключает} \ru{каналы}.
\end{quote}

\begin{quote}
\emph{Pápa vsyo vrémya pereklyucháyet kanály.}
\end{quote}

\begin{quote}
Dad keeps switching channels.
\end{quote}

Imperfective; perfective \ru{переключить}. \emph{\ru{Переключать} \ru{на} \ru{другой} \ru{канал}}.

\begin{cousinweb}
\emph{\ru{переключать}} is built on \emph{\ru{пере}-} and the root of \emph{\ru{ключ}} 'a key, a switch'; EtymDB lists the verb without a source.
\end{cousinweb}

\begin{grammarlens}[title={What you've built}]
One more word for telling, explaining and describing at B1.
\end{grammarlens}

\subsection*{Your turn}
\begin{itemize}
\raggedright
  \item \emph{Say it:} \emph{pereklyuchát}
  \item \emph{Say it:} \emph{pereklyuchát}, once more
  \item \emph{Say it:} a sentence of your own with \emph{pereklyuchát}
  \item \emph{Recall:} say the Russian for an episode, then the Russian for subtitles, then say \emph{pereklyuchát} again
\end{itemize}

\begin{tcolorbox}[breakable,colback=teal!4,colframe=teal!35!black,title={Before you move on}]
What does \ru{переключать} mean? (\textquotedblleft{}To switch (channels)\textquotedblright{}.) Say it once more.
\end{tcolorbox}
\section[\texorpdfstring{\ru{экранизация} (ekranizátsiya)}{экранизация (ekranizátsiya)}]{\ru{экранизация} (ekranizátsiya) — a screen adaptation}

\label{lesson:RU-C527-ekranizatsiya}



\begin{quote}
Before the new one: say the Russian for subtitles, then the Russian for to switch.
\end{quote}

\subsection*{You'll want to know: \ru{экранизация}}
\textbf{\ru{экранизация}} (\emph{ekranizátsiya}) — \textquotedblleft{}a screen adaptation\textquotedblright{}.

\begin{quote}
\ru{Экранизация} \ru{оказалась} \ru{лучше} \ru{книги}.
\end{quote}

\begin{quote}
\emph{Ekranizátsiya okazálas lúchshe knígi.}
\end{quote}

\begin{quote}
The film version turned out better than the book.
\end{quote}

Feminine.

\begin{cousinweb}
\emph{\ru{экранизация}} is built on \emph{\ru{экран}} 'a screen', which was borrowed from French \emph{écran} (EtymDB); EtymDB has no entry for \emph{\ru{экранизация}}.
\end{cousinweb}

\begin{grammarlens}[title={What you've built}]
One more word for telling, explaining and describing at B1.
\end{grammarlens}

\subsection*{Your turn}
\begin{itemize}
\raggedright
  \item \emph{Say it:} \emph{ekranizátsiya}
  \item \emph{Say it:} \emph{ekranizátsiya}, once more
  \item \emph{Say it:} a sentence of your own with \emph{ekranizátsiya}
  \item \emph{Recall:} say the Russian for subtitles, then the Russian for to switch, then say \emph{ekranizátsiya} again
\end{itemize}

\begin{tcolorbox}[breakable,colback=teal!4,colframe=teal!35!black,title={Before you move on}]
What does \ru{экранизация} mean? (\textquotedblleft{}A screen adaptation\textquotedblright{}.) Say it once more.
\end{tcolorbox}
